\section{The zero-spare binary machine}
\label{sec:z215-machine}

The binary construction changes the policy for allocating unused register
blocks.  Its source program is the accepted $4{,}751$-instruction K215 source;
the logical operations on its nine primary registers remain the same.  The
new initial state reaches a different source seed, whose soundness requires
the two-truth-prefix theorem proved below.

\subsection{Exact table and state conventions}

Use one two-way infinite tape, alphabet $\{0,1\}$, an initially zero tape,
initial head position $0$, and moves $L,R$ only.  Native labels are retained
from the accepted K215 table.  Make exactly these changes:
\[
 \text{initial state: }0\longmapsto74,
 \qquad
 \delta(198,0): (1,L,40)\longmapsto(0,L,40).
\]
Every other transition, including both rows of state $0$ and the row
$\delta(198,1)=(1,L,113)$, is unchanged.  The resulting table has
$215$ working states and $430$ rows.  Native halt label $254$ has no outgoing
rows and is excluded from the count.  All $215$ working states are reachable
from $74$ in the finite transition graph; this graph assertion makes no claim
that all states occur on the particular blank run.

For the dense presentation, list native state $74$ first and the remaining
working labels in increasing order, then number this list $0,\ldots,214$.
Map native halt $254$ to dense halt $215$.  Native dispatcher $84$ becomes
dense dispatcher $83$.  The supplied canonicalization file records both
directions of this exact bijection.  Thus the delivered dense machine starts
in state $0$ and has the usual consecutive working labels.

This construction is not a state renaming of K215.  The old table has $189$
rows writing symbol $1$; the new table has $188$.  That count is invariant
under state renaming.  More substantively, every source increment now appends
a zero-valued spare rather than a one-valued spare, changing tape growth
throughout an unbounded run.

\subsection{Physical frames and the spare policy}

Write
\[
 B(v)=0\,1^{v+1},\qquad C(v)=1\,0\,1^v \quad(v\in\mathbb N).
\]
The primary logical register order is
$(N,T,W,A,D,H,X,Y,L)$.  The physical file order is
$(W,L,T,D,A,Y,N,H,X)$, followed by a nonempty finite list $S$ of natural-valued
spares.  No theoremhood condition is imposed on these spare values.

A frame has the dispatcher at head position $h$, a marker $1$ at $h-1$, a
binary control address $p$ starting at $h$, enough following zeros to put a
header $1$ at $h+64+g$, and then the nine primary blocks and the spare blocks.
Here $g\in\mathbb N$, $|p|\leq64$, and the whole remaining exterior is zero.
In symbols, its finite distinguished part is
\[
 1\;p\;0^{64+g-|p|}\;1\;
 B(W)B(L)B(T)B(D)B(A)B(Y)B(N)B(H)B(X)B(S).
\]
The head is at the beginning of $p$; if $p$ is empty it scans the first zero
after the marker.  The notation $B(S)$ means concatenation of the spare
blocks.  Addresses with trailing zeros are interpreted by the unchanged
binary source dispatcher.  In particular the empty address denotes the
source root.

Let the frame projection forget $h,g,S$ and retain the resolved source
instruction together with the nine logical registers.  The change at198/0
affects only the newly allocated spare.  Under the accepted engine it is
$B(1)=011$; here it is $B(0)=01$.  Each corresponding increment frame is
therefore one tape cell shorter at the end of its file.  No primary value or
source successor is changed by this allocation choice.

Starting from otherwise matched frames, after $k$ source increments the
newly appended spare region has $2k$ cells and $k$ ones, instead of
$3k$ cells and $2k$ ones. Thus that region is exactly one third shorter.
This comparison is at corresponding source boundaries; it is not a claim
about total tape space at equal binary times or about runtime speed.

\begin{lemma}[Quantified primitive simulation]
\label{lem:z215-primitive}
For every primitive occurrence of the pinned source, every applicable
decrement guard, every natural primary register tuple, every natural gap,
and every nonempty finite natural spare list, the corresponding Z215 frame
executes a finite positive number of binary transitions to its prescribed
successor frame.  An increment appends one zero spare.  Other nonhalting
primitives preserve the spare list.  A source halt reaches the binary halt;
no nonhalting primitive block reaches it.
\end{lemma}

\begin{proof}
A fresh execution of the supplied symbolic prover derives all
$6{,}176$ guarded cases for the $4{,}751$ primitive occurrences. The saved
per-case records are receipts of that execution, not a dump of every
derivation event; the complete process is reproducible with the included
code. It uses the unchanged inherited checker, with its existing allocation
parameter set to zero. The argument represents complete tapes by finite symbolic words,
unary runs of affine length, finite lists of blocks, and zero exteriors.

Its elementary rules are literal transition steps, induction over a unary
run when the exact row is a scan-direction self-loop, and induction over a
finite block list after a one-block lemma has been derived from the exact
table.  The four block lemmas have native entries $64,0,40,97$: respectively
rightward $B$ preservation, rightward $B$-to-$C$ conversion, leftward $B$
preservation, and leftward $C$-to-$B$ conversion.  Each crosses one arbitrary
block of value $v$ in exactly $v+2$ steps.  Concatenation uses the elementary
identity $C(S)1=1B(S)$ and its reversal.  These are unary and list inductions,
not tests on a bounded set of register values.

Thirteen affine slots represent a constant, the gap, nine independent
primary parameters, the first spare value, and the total width of the
remaining spare list.  A positive decrement uses one plus its parameter
for the selected register; its zero branch fixes that register to zero.
Every resulting duration has nonnegative integer coefficients and a positive
constant.  The minimum constant across the complete cases is $6$.

The derivations contain $1{,}698{,}790$ symbolic rule applications.  Together
with the four block schemas they account for all $430$ literal rows.  The
ordered case file binds each result to its source occurrence and guard;
the exact table and source hashes bind the certificates to the delivered
bytes.  Every case closes with the specified whole-tape endpoint, including
the unbounded zero exterior, proving the stated contract.
\end{proof}

For clarity, the exact endpoint also tracks counter carries and head motion.
If $\varepsilon\in\{0,1\}$ is the source address-increment overflow and
$d=1$ for a positive decrement, $d=0$ otherwise, then the successor frame
has head $h-\varepsilon$ and gap $g+d+\varepsilon$.  Continue instructions
have no increment overflow.  These changes disappear under the frame
projection; they are retained in the whole-tape certificates.

\subsection{The complete blank prefix}

Starting at native $74$ on the zero tape, the machine reaches the source-root
frame after exactly $117{,}004$ transitions.  Its values are
\[
 \begin{gathered}
 q=84,\qquad h=-4,\qquad \text{header}=80,\qquad g=20,\\
 (N,T,W,A,D,H,X,Y,L)=(1,0,5,0,0,0,0,0,0).
 \end{gathered}
\]
There are exactly $67$ spare blocks.  The third and eighth have value $1$;
every other spare has value $0$.  The physical primary values are
$(5,0,0,0,0,0,1,0,0)$.  The complete tape has $86$ one-cells, with least
position $-5$ and greatest position $240$; all other cells are zero.

These assertions are an exact finite computation.  An independent audit
replays every one of the $117{,}004$ transitions in both the native and dense
presentations, then reconstructs the entire terminal frame from its stated
primary and spare blocks.  The reconstructed set of one-cells equals the
replayed set.  Thus the initial configuration of the source theorem is
established, including the otherwise essential dirty stack value $W=5$;
a merely plausible-looking intermediate register frame is not being used.

\begin{theorem}[Frame projection and positive-duration lift]
\label{thm:z215-lift}
From the displayed blank-start frame, the projection of Z215 at successive
primitive boundaries is exactly the run of the pinned source from
$(1,0,5,0,0,0,0,0,0)$.  The binary machine halts if and only if that source
run halts.
\end{theorem}

\begin{proof}
Apply Lemma~\ref{lem:z215-primitive} inductively.  The successor spare list
remains finite and nonempty, so the lemma can always be applied again.  It
preserves the exact logical effect and source successor, while the projection
forgets harmless changes of position, gap, and spare contents.  Every
nonhalting primitive has a finite duration at least $6$.  Consequently an
infinite source run yields infinitely many binary transitions, and these
finite consecutive blocks exhaust its deterministic binary execution.
There can be no extra binary halt between source boundaries.  Conversely,
a finite source run ending at its halt instruction reaches the binary halt
in that instruction's finite block.  Prepending the exact blank prefix
establishes the equivalence from the actual zero tape.
\end{proof}

Combining this lift with the two-truth-prefix source theorem and inherited
numerical-calculus calibration gives, by ordinary reasoning formalizable in
ZFC,
\[
 \mathrm{ZFC}\ \vdash\
 \bigl(\operatorname{Halt}(\mathrm{Z215},\mathrm{blank})
       \leftrightarrow\neg\operatorname{Con}(\mathrm{ZFC})\bigr).
\]
Equivalently its blank nonhalting is provably equivalent to
$\operatorname{Con}(\mathrm{ZFC})$.  The finite symbolic certificates and
their induction rules supply an ordinary mathematical proof.  They are not
a newly compiled exact-table Lean or Rocq theorem, nor a new internal-PA
endpoint.  No unconditional nonhalting or state-minimality claim is made.
